\section{Moving centers and scalar vanishing}
\label{sec:moving-centers}

We work on the terminal fourfold $Y$ from \Cref{sec:reduction}, write
$K=K_Y$, and put $n=4$. Write $Y^{\mathrm{vg}}$ for a fixed very
general locus on which both \Cref{prop:reduction,prop:subvarieties}
apply. Thus every curve through this locus has $K$-degree at least
$1/\iota$, and every proper positive-dimensional subvariety through
it is of general type. We will bound the order of every section of
every Cartier multiple of a suitable ample rational divisor at a
very general point. Fix a very ample Cartier divisor $A$ on $Y$.
For a fixed positive rational $\epsilon$ and positive rational $t$
tending to zero, define
\begin{equation}
\label{eq:small-volume-polarization}
 \mu_t=\bigl((K+tA)^n\bigr)^{1/n},\qquad
 r_t\in\Z_{>0},\quad r_t\mu_t\longrightarrow\epsilon,
 \qquad L=L_t=r_t(K+tA).
\end{equation}
Such integers $r_t$ exist, for example by rounding $\epsilon/\mu_t$.
Since $K$ is nef and not big, $\mu_t>0$ tends to zero. Hence
\begin{equation}
\label{eq:small-volume-curve-lower-bound}
 r_t\longrightarrow\infty,\qquad
 L_t^n\longrightarrow\epsilon^n,
 \qquad L_t\cdot C\geq r_t/\iota
\end{equation}
for every integral curve $C$ through that very general locus.
The integers $r_t$ are unrelated to the original Cartier index $r$
in \Cref{cor:lc-nonvanishing}.

A section with excessive vanishing will yield a curve through
$Y^{\mathrm{vg}}$ whose $L_t$-degree is bounded independently of $t$,
contradicting \Cref{eq:small-volume-curve-lower-bound}. To construct
that curve, we follow the base loci of the full jet kernels as their
marked point moves. A component that persists at two successive
jet levels has high multiplicity; this bounds both its degree and
its canonical intersection. General type then produces a curve
of bounded degree. We first give the three moving-center estimates
from \cite[Lemmas~7.1--7.4]{OpenAIAbundance2026}, including the
differentiation argument responsible for ordinary ideal powers.
These inputs are independent of abundance and subadditivity. We
then prove the curve obstruction and apply it to scalar vanishing.

\subsection{Numerical subadjunction at a generic lc center}

The numerical statement below only requires log canonicity near
the generic point of the center. This is the form needed for
boundaries obtained from a linear system.

The subadjunction method relates an lc center to the canonical bundle
formula. Kawamata's minimal-center theorem and the later unperturbed
form of Fujino--Gongyo are important antecedents
\cite{KawamataSubadjunction1998,FujinoGongyoSubadjunction2012}.
The numerical form below permits a center that is only generically lc
and need not be minimal. Its proof in the companion uses a dlt stratum,
the generic rank-one condition, b-nef moduli, and finite-base ramification;
it is not an application of those antecedents with their hypotheses
omitted.

\begin{lemma}[Numerical generic subadjunction]
\label[lemma]{lem:numerical-subadjunction}
Let $Z$ be a projective $\Q$-factorial klt variety over $\C$, let
$\Theta\geq0$ be a rational divisor, and let $V\subset Z$ be an
integral positive-dimensional subvariety. Suppose that $(Z,\Theta)$
is lc at the generic point of $V$ and that a divisor of log
discrepancy zero has center $V$. Let $f=\dim V$, let $P$ be a nef
rational Cartier divisor on $Z$, and let $\rho:V^*\to V$ be a smooth
projective resolution, factoring through the normalization. Then
\begin{equation}
\label{eq:numerical-subadjunction}
 K_{V^*}\cdot(\rho^*P)^{f-1}
 \leq (K_Z+\Theta)|_V\cdot(P|_V)^{f-1}.
\end{equation}
The right side is the intersection of the restricted rational
Cartier class with the cycle $[V]$. No $\Q$-Gorenstein hypothesis
on the normalization of $V$ is required.
\end{lemma}

\noindent The proof is given in the independent moving-center part of
\cite[Lemma~7.1]{OpenAIAbundance2026}.\par


\subsection{Degree and canonical bounds for a base component}

The following companion lemma applies numerical subadjunction to a boundary built from an
actual linear system. Its local multiplicity controls both the
degree of the center and the coefficient of that boundary.

\begin{lemma}[A base component of high multiplicity]
\label[lemma]{lem:base-component}
Let $Z$ be a projective $\Q$-factorial klt variety of dimension $d$,
let $P$ be an ample rational Cartier divisor, and let $k>0$ be an
integer for which $kP$ is Cartier. Let
$0\ne\mathcal H\subset H^0(Z,\cO_Z(kP))$ be a linear subspace
with proper base locus and base ideal $\mathfrak b$.
Suppose that $V$ is an integral base-locus component of dimension
$f<d$, isolated at its generic point $\eta_V$, and that
$\eta_V\in Z_{\rm sm}$. Write $a=d-f$ and
$\mathfrak m=\mathfrak m_{\cO_{Z,\eta_V}}$. If
\[
 \mathfrak b\cO_{Z,\eta_V}\subset\mathfrak m^h
 \qquad(h\in\Z_{>0}),
\]
then
\begin{equation}
\label{eq:base-component-degree}
 P^f\cdot V\leq (k/h)^aP^d.
\end{equation}
If $f>0$, there is a rational number $c$ with
$0<c\leq ak/h$ such that, on a smooth resolution $V^*$,
\begin{equation}
\label{eq:base-component-canonical}
 K_{V^*}\cdot P^{f-1}
 \leq (K_Z+cP)|_V\cdot P^{f-1}.
\end{equation}
The coefficient $c$ may be chosen with an effective rational divisor
$\Theta\sim_{\Q}cP$ such that $(Z,\Theta)$ is lc at the generic
point of $V$ and has an lc place with center $V$. Thus
\Cref{lem:numerical-subadjunction} can also be applied with a different
nef testing class. The degree estimate includes $f=0$; the canonical estimate is
asserted only for positive-dimensional $V$.
\end{lemma}

\noindent The proof is given in the independent moving-center part of
\cite[Lemma~7.2]{OpenAIAbundance2026}.\par


\subsection{Differentiating the full moving jet kernel}

We use the parameter-differentiation method developed in
\cite[Proposition~2.3]{EinKuchleLazarsfeld1995}. The form used here
retains ordinary ideal powers and isolated base components. Fixing a
high order at a varying point produces a family of
linear subspaces of one fixed section space. A component common
to two successive base loci acquires high multiplicity, because
parameter derivatives of the higher kernel still belong to the
lower kernel. We state both the parameter-family conclusion and
the local slicing property that will be used later.

\begin{lemma}[Moving multiplicity]
\label[lemma]{lem:moving-multiplicity}
Let $Z$ be an integral projective variety of dimension $d$, and
let $P$ be an ample rational Cartier divisor. Fix an integer $k$
with $kP$ Cartier and positive real numbers
\[
 \tau_0<\tau_1<\cdots<\tau_d,
 \qquad\tau_{i+1}-\tau_i=\delta>0.
\]
For a smooth point $x$, let $\mathcal H_i(x)$ be the entire
subspace of sections of $kP$ vanishing at $x$ to order at least
$\ceil{k\tau_i}$. Suppose that for very general $x$ all these
subspaces are nonzero and the base locus of $\mathcal H_0(x)$
has a positive-dimensional component through $x$.
If $k\delta\geq4$, then, after an algebraic parameter extension
and restriction to nonempty opens, there are an irreducible
parameter space $T$, a fixed adjacent pair of levels, and a
family of integral subvarieties $V_t$ common to the two base
loci such that the incidence
\[
 \Gamma=\{(t,z):z\in V_t\}\subset T\times Z
\]
dominates $Z$. Each general $V_t$ has dimension in $[1,d-1]$,
is an isolated component of the higher base locus at its generic
point, and its entire higher-series base ideal is contained
there in
\begin{equation}
\label{eq:moving-multiplicity-order}
 \cI_{V_t}^{h},\qquad h=\ceil{k\delta/2},
 \qquad k/h\leq2/\delta.
\end{equation}
These are ordinary local ideal powers in the smooth ambient open.

More precisely, the incidence base ideal lies in $\cI_\Gamma^h$
on a dense open of $\Gamma$. If a further morphism $Z\to B$ is
given, one may choose a general incidence point where the base
support is only $V_t$ and $V_t$ is smooth over the smooth open of
its image in $B$. For the fiber component $F$ through that point,
assume that $F$ is proper in the integral ambient fiber component
under consideration. The restricted series is then nonzero and
has base ideal in the ordinary
$\cI_F^h$ and has $F$ as an isolated base component there.
This last assertion concerns the scheme fiber locally at that
chosen point, and does not assert transversality of every fiber.
Smoothness near the generic point of $F$, and the projectivity and klt
hypotheses needed for \Cref{lem:base-component} on the ambient fiber,
must be checked in each application.
\end{lemma}

\begin{proof}
We recall the argument of \cite[Lemmas~7.3--7.4]{OpenAIAbundance2026}
to explain why the estimate concerns the entire kernel and ordinary
ideal powers. On the open where the jet maps have constant rank,
their kernels are vector bundles with fibers $\mathcal H_i(x)$.
Their base loci increase with $i$. Starting with a positive-dimensional
component through the mark, follow a containing component at each
level. The $d+1$ components all have dimensions in $[1,d-1]$, so
two successive ones coincide. After a finite parameter extension
and shrinking, their geometric generic components and inclusions
spread to one family $\Gamma\subset T\times Z$. Its evaluation
dominates $Z$, since each member contains its moving mark $x(t)$.

Take a local frame of the higher kernel and trivialize $kP$ near
$z$ with a frame independent of $t$. In a fixed basis of the global
section space, each frame section has the form
\[
 s(t,z)=\sum_\alpha c_\alpha(t)s_\alpha(z).
\]
Differentiation in the parameter with $z$ fixed therefore gives
another global section of $kP$. A derivative of order $r$ lowers
the vanishing order at $x(t)$ by at most $r$. It belongs to the
entire lower kernel whenever
\[
 r\leq\ceil{k\tau_{i+1}}-\ceil{k\tau_i}.
\]
The right side is at least $k\delta-1$, so every derivative of
order less than $h=\ceil{k\delta/2}$ belongs to that kernel and
vanishes on $\Gamma$. This is where using the full kernels matters:
an arbitrarily chosen subseries need not contain these derivatives.

To turn these vanishings into ordinary ideal membership, work at a
general smooth incidence point. Since $\Gamma\to Z$ is submersive,
parameter directions span its normal space in $T\times Z$.
Choose parameter coordinates $t=(t',t'')$ so that the implicit
function theorem writes $\Gamma$ as $t'=F(t'',z)$. With $t''$ and
$z$ fixed, derivatives in $t'$ are exactly derivatives in the
normal coordinates $u=t'-F(t'',z)$. All normal Taylor coefficients
of degree below $h$ therefore vanish, proving
$s\in\cI_\Gamma^h$. Faithful flatness of analytic local rings over
algebraic local rings gives the same membership algebraically.
There are finitely many frame sections, so this containment holds
for the entire incidence base ideal on a dense algebraic open.

Generic smoothness of $\Gamma\to T$ identifies its scheme fiber
there with the reduced member $V_t$. Restriction gives
$\cI_\Gamma^h\cO_{\{t\}\times Z}=\cI_{V_t}^h$, and the frame
specializes to a basis of the entire higher kernel. Isolation
holds because $V_t$ is a base-locus component. Finally remove all
other base components and choose the incidence point where
$V_t$ is smooth over its actual image in $B$. Its scheme fiber is
then reduced locally, so restriction to the ambient fiber sends
$\cI_{V_t}^h$ to $\cI_F^h$ and the base support to $F$. If $F$
is proper in the chosen integral ambient fiber component, that
support is proper there; hence the restricted series is nonzero
and has $F$ as an isolated base component at its generic point.
\end{proof}


For a fixed gap $\delta$, combining
\Cref{lem:moving-multiplicity,lem:base-component} gives degree and
canonical-intersection bounds for a repeated moving component,
with coefficients independent of the large divisible integer $k$.
The next lemma converts such intersection bounds into a curve of
bounded degree. Its general-type hypothesis supplies a birational
linear system in a degree depending only on the dimension. The
scalar application uses \Cref{prop:subvarieties} and has no boundary.
We include a reduced boundary in the statement because the
two-slot argument will also need its logarithmic form.

\subsection{A uniform obstruction from logarithmic general type}

\begin{lemma}[Bounded-degree curve obstruction]
\label[lemma]{lem:degree-obstruction}
Keep $Y$, $L_t$, and $\iota$ as in
\Cref{eq:small-volume-polarization,eq:small-volume-curve-lower-bound}.
Fix $f\geq1$ and positive constants $C_0,C_1$. Suppose, for
arbitrarily small $t$, that there are a smooth projective variety
$V_t^*$ of dimension $f$, a reduced simple normal crossings
divisor $G_t$ on it, a nef rational divisor $P_t$, and a morphism
$e_t:V_t^*\to Y$ satisfying the following conditions:
\begin{enumerate}[label=\textup{(\roman*)}]
\item $K_{V_t^*}+G_t$ is big;
\item $0<P_t^f\leq C_0$ and
$(K_{V_t^*}+G_t)\cdot P_t^{f-1}\leq C_1P_t^f$;
\item $e_t$ is generically finite onto its positive-dimensional
image, and that image meets the fixed very general locus of $Y$;
\item $P_t-e_t^*L_t$ is nef.
\end{enumerate}
Then these conditions are impossible for all sufficiently small
$t$. The case $G_t=0$ includes ordinary general type.
\end{lemma}

\begin{proof}
Uniform effective birationality for projective lc pairs of
dimension $f$ with big adjoint and coefficients in the fixed DCC
set $\{0,1\}$ supplies an integer $b=b(f)$ such that
$|b(K_{V_t^*}+G_t)|$ is birational
\cite[Theorem~4.0.1]{HMX2018}. Resolve its base ideal and denote
the free moving divisor by $H_b$. We keep the same notation for
the pullbacks of $P_t$ and $L_t$. The fixed divisor is effective,
so
\[
 H_b\cdot P_t^{f-1}
 \leq b(K_{V_t^*}+G_t)\cdot P_t^{f-1}
 \leq bC_1P_t^f.
\]
Both $H_b$ and $P_t$ are nef and big; $P_t$ is big because
$P_t^f>0$. The Khovanskii--Teissier log-concavity inequalities for
their mixed intersections give, when $f>1$,
\begin{equation}
\label{eq:degree-obstruction-mixed}
 P_t\cdot H_b^{f-1}
 \leq
 \frac{(P_t^{f-1}\cdot H_b)^{f-1}}{(P_t^f)^{f-2}}
 \leq(bC_1)^{f-1}P_t^f
 \leq(bC_1)^{f-1}C_0.
\end{equation}
Indeed the successive ratios of the positive numbers
$P_t^{f-i}H_b^i$ decrease, which gives the first inequality.

Choose a point where the birational morphism defined by $H_b$
is an isomorphism onto an open of its image, where $e_t$ is
quasi-finite, and whose image lies in the very general locus of
$Y$. These conditions are compatible. To spell out the last
one, write the excluded subset of $Y$ as a countable union of
proper closed subsets. The image of $e_t$ is not contained in
any of them because it meets their complement. Their inverse
images are therefore proper closed subsets of $V_t^*$. They
can be avoided together with the finitely many dense-open
conditions. For a sweeping family one first chooses a very
general member not contained in any excluded subset, and then
such a point on that member.

For $f>1$, cut by $f-1$ general members of $|H_b|$ passing through
the chosen point. These give a proper effective curve cycle:
on the birational image they are general hyperplanes through a
point of the isomorphism open, and their pullbacks have no
positive-dimensional base component. A component $C$ through
the selected point maps nonconstantly to $Y$, since $e_t$ is
quasi-finite there. Nefness and
\Cref{eq:degree-obstruction-mixed} give
\[
 L_t\cdot e_t(C)
 \leq e_t^*L_t\cdot C
 \leq P_t\cdot C
 \leq P_t\cdot H_b^{f-1}
 \leq(bC_1)^{f-1}C_0.
\]
In the first inequality the positive degree of the finite map
of the normalization of $C$ onto its image has merely been
discarded. If $f=1$, use $V_t^*$ itself as $C$; the same argument
gives $L_t\cdot e_t(C)\leq P_t\cdot C\leq C_0$.

The image is an integral curve through the very general locus,
so \Cref{eq:small-volume-curve-lower-bound} bounds its degree
below by $r_t/\iota$. This tends to infinity, whereas the upper
bound depends only on $f,C_0,C_1$. The contradiction proves the
lemma. If several dimensions $1\leq f\leq n$ occur, take the
maximum of the finitely many resulting constants.
\end{proof}

\subsection{The scalar order bound}

Define
\begin{equation}\label{eq:scalar-polarization}
 P_{\mathrm s}=2r_t(K+2tA),\qquad
 \rho_t=(P_{\mathrm s}^{\,n})^{1/n}.
\end{equation}
The inequalities
\[
 2L\le P_{\mathrm s}\le4L
\]
are inequalities in nef order. After decreasing \(t\), they imply
\begin{equation}\label{eq:scalar-volume-bounds}
 \epsilon\le\rho_t\le8\epsilon .
\end{equation}
Only the existence of a positive lower bound and the displayed upper
bound will be used.

\begin{proposition}[Scalar orders]\label[proposition]{prop:scalar-orders}
For all sufficiently small positive rational \(t\), a very general
point \(x\in Y_{\mathrm{sm}}\) has the following property. For every
positive integer \(k\) for which \(kP_{\mathrm s}\) is Cartier and every
nonzero section
\(s\in H^0(Y,\cO_Y(kP_{\mathrm s}))\), one has
\begin{equation}\label{eq:scalar-orders}
 \ord_x(s)\le4k\rho_t\le32k\epsilon .
\end{equation}
\end{proposition}

\begin{proof}
Fix \(t\) for the moment. For each Cartier multiple \(kP_{\mathrm s}\)
and each integer \(a\), the condition
\[
 H^0(Y,\cO_Y(kP_{\mathrm s})\otimes\cI_x^a)\ne0
\]
is a rank condition on the jet evaluation map over \(Y_{\mathrm{sm}}\).
We choose \(x\) outside all proper rank-jumping loci of these
countably many maps. Thus if the order bound fails at such a point,
the relevant high-order kernel is nonzero for general marks as well.
Taking powers of the offending sections gives the same strict
inequality in arbitrarily large divisible degrees. This permits all
subsequent choices of a sufficiently large divisible \(k\).

Choose \(n+1\) levels equally spaced from \(2\rho_t\) to \(3\rho_t\);
their gap is
\[
 \delta_{\mathrm s}=\rho_t/n .
\]
If \(\ord_x(s)>4k\rho_t\) in a sufficiently large degree, every one of
these jet kernels is nonzero. The base locus of the lowest kernel
cannot have \(x\) as an isolated component. Indeed its base ideal at
\(x\) is contained in
\(\mathfrak m_x^{\ceil{2k\rho_t}}\), so the zero-dimensional case of
\Cref{lem:base-component} would give
\[
 1\le
 \left(\frac{k}{\ceil{2k\rho_t}}\right)^nP_{\mathrm s}^{\,n}
 \le2^{-n},
\]
a contradiction.

Apply \Cref{lem:moving-multiplicity}. It produces a dominant family
of proper positive-dimensional base components \(V\), of some
dimension \(1\le f<n\), such that the entire higher jet kernel's generic
base ideal along \(V\) is primary and contained in
\(\cI_V^h\), where
\[
 h=\ceil{k\delta_{\mathrm s}/2},
 \qquad k/h\le2/\delta_{\mathrm s}.
\]
For a general member, \Cref{lem:base-component} gives
\begin{align}
 P_{\mathrm s}^{\,f}\cdot V
 &\le(2/\delta_{\mathrm s})^{n-f}P_{\mathrm s}^{\,n},
 \label{eq:scalar-center-volume}\\
 K_{V^*}P_{\mathrm s}^{\,f-1}
 &\le
 \left(\frac1{2r_t}+\frac{2(n-f)}{\delta_{\mathrm s}}\right)
 P_{\mathrm s}^{\,f}\cdot V ,
 \label{eq:scalar-center-canonical}
\end{align}
where \(V^*\) is a smooth projective resolution. The second inequality
uses \(K\le P_{\mathrm s}/(2r_t)\).

The constants on the right are bounded independently of small \(t\),
of \(k\), and of the member \(V\), by
\Cref{eq:scalar-volume-bounds}. We choose a very general incidence
point before choosing the member \(V\). Its image belongs to
\(Y^{\mathrm{vg}}\), so \(V^*\) is of general type by
\Cref{prop:subvarieties}. Moreover \(P_{\mathrm s}\) dominates \(L\)
in nef order. The inclusion of \(V\) into \(Y\), composed with its
resolution, satisfies all the hypotheses of
\Cref{lem:degree-obstruction}. \Crefrange{eq:scalar-center-volume}{eq:scalar-center-canonical}
would therefore give a curve through \(Y^{\mathrm{vg}}\) of bounded
\(L\)-degree, contradicting
\Cref{eq:small-volume-curve-lower-bound} as \(t\) becomes small.

The bounds used in this contradiction do not depend on the initial
degree in which excessive vanishing occurred. Consequently, after
one restriction to sufficiently small \(t\), all Cartier multiples
satisfy \Cref{eq:scalar-orders} at a very general point.
\end{proof}
